%%%%%%%%%%  scan idx 264  |  folio 253  %%%%%%%%%%

\medskip
\noindent 2.4. Sea
\[
 f:\underline T'\longrightarrow\underline T
\]
un morfismo de topoi. Se tiene entonces un homomorfismo canónico
(VII 2.8)
\[
 \Biext^1(P,Q;G)\longrightarrow\Biext^1(P',Q';G'),
\]
donde las primas denotan los grupos imagen inversa. Como para todo
grupo conmutativo $E$ de $\underline T$ existe manifiestamente un
isomorfismo canónico
\[
 ({}_nE)'\simeq{}_n(E'),
\]
los homomorfismos (2.2.2) sobre $\underline T$ y $\underline T'$ dan
un diagrama canónico
\begin{equation}\tag{2.4.1}
\begin{tikzcd}[column sep=huge,row sep=huge]
 \Biext^1(P,Q;G)\arrow[r]\arrow[d,"{\varphi^n_{P,Q}}"']
   &\Biext^1(P',Q';G')\arrow[d,"{\varphi^n_{P',Q'}}"]\\
 \Hom({}_nP\otimes{}_nQ,G)\arrow[r]
   &\Hom({}_nP'\otimes{}_nQ',G').
\end{tikzcd}
\end{equation}
De manera análoga, los homomorfismos (2.2.3) dan un diagrama semejante
en términos de los $T_\ell$. Afirmo que estos diagramas son
\underline{conmutativos}. La verificación es esencialmente inmediata
a partir de las definiciones y se deja al lector. Más generalmente,
obsérvese que \underline{todos} los homomorfismos canónicos
considerados en esta sección y en la precedente son
\underline{compatibles con el cambio de topos}, en el sentido
evidente. En adelante utilizaremos sin mencionarlas las compatibilidades
de este tipo.
